\section{Introduction}
\label{sec:intro}

Learned world models have been studied primarily in settings where
the environment is observed as pixels and the task is to predict the
next frame or latent video token. The standard recipe pairs a
high-capacity visual encoder--decoder---typically a CNN, U-Net, or
video Transformer---with an action-conditioned dynamics module. This
is a natural fit when the observation is visually rich and the
underlying mechanics are comparatively smooth or uniform. It is a
less natural fit for games whose difficulty lies not in rendering but
in a large space of symbolic rule systems.

PuzzleScript games are such a regime. Each game is fully specified by
a short text file describing object kinds, layered grid levels, and a
list of declarative rewrite rules, e.g.
\texttt{[ > Player | Crate ] -> [ > Player | > Crate ]}. The engine
evaluates these rules in order on the grid at every tick, applying each rule until the state converges, with \texttt{again} rules causing the entire rule block to fire again. The state can be represented as a discrete multi-channel
grid, the dynamics are deterministic and combinatorial, and
the hard computation per tick is a sequence of atomic local updates
applied repeatedly.

To capture this symbolic mechanical complexity, we propose a \emph{rule-conditioned Neural
Cellular Automaton} (NCA) world model.
A perceiver-style encoder maps
the tokenized game description to $K$ latent slots; the dynamics body
is a local update rule applied for $T$ iterations on a hidden grid,
with cells cross-attending to those slots at every step. The same body
parameters can be reused across iterations, mirroring the engine's
outer loop, while the conditioner makes the dynamics model portable
across games. In the full version of the model, these slots are tied
to a symbolic-game autoencoder so that game latents can both condition
dynamics and decode back into program text.
Foregoing pixels lets us model the environment in
its native state representation and measure fidelity cell-by-cell
against the original engine.

Our corpus contains 3{,}474 token-distinct
human-authored games (one to twenty rules each), of which up to 199
are used in the largest training preset. Across these games we cache
bounded-A* transitions from authored levels, yielding $\sim$15M
transitions; held-out evaluation uses a 30-game stratified sample
disjoint from training by name and token-sequence hash.

On the 94-game broad-corpus training preset, the trained NCA achieves
a median per-cell teacher-forced one-step error of $0.15\%$ and a
median 50-step random-action autoregressive rollout error of
$2.5\%$ across all 94 games and all authored levels.
% We use ``near-perfect'' to mean a one-step teacher-forced cell-error rate below $1\%$ on the broad corpus, with the multi-step rollout numbers quantifying the drift caused by autoregressive compounding.
To rule out per-state memorization, we also report a
within-game OOD-levels probe (\S\ref{sec:results-within-game}) in
which every authored level is held out from training and only
purely-synthetic rollouts at the same grid size enter the training
set; the same rule-conditioned NCA recovers near-zero rollout error
on all authored levels of Sokoban-class games and non-trivial
transfer on three further mechanic classes, evidencing that the
high-precision regime reflects learned rule structure rather than
authored-state lookup.

\paragraph{Contributions.}
\begin{enumerate}
  \item \textbf{NCA dynamics over discrete state.}
        Modeling the engine's native discrete grid with a shared,
        iteratively-applied local update body matches the reference
        dynamics with median 1-step teacher-forced cell-error of
        $0.15\%$ across the 94-game broad-corpus preset. A controlled
        comparison that holds the rule encoder fixed and varies only
        the spatial body shows the iterated-local-update bias
        outperforms non-iterative architectures (CNN, U-Net, ViT) on
        rollout fidelity at matched or smaller parameter counts.
        Within-game OOD-levels probes (\S\ref{sec:results-within-game})
        confirm the model is learning rule structure rather than
        memorizing authored states.
  \item \textbf{Rule-conditioning for multi-game simulation.}
        A single model simulates dozens of mechanically diverse games
        by cross-attending to slot embeddings produced by an encoder
        over each game's tokenized rule text. Trained jointly with a
        slot-token decoder, the same latent drives forward dynamics
        and decodes back to program text: $76/94$ training games
        reconstruct perfectly under teacher forcing, $15/25$ samples
        from the empirical slot prior decode to engine-loadable
        PuzzleScript, and slot-space interpolation between two games
        is engine-loadable at $9/9$ points along the line.
        The same latent is also recoverable by inverse-fitting through
        the frozen NCA, driving per-cell BCE 2--3 orders of magnitude
        down on held-out targets and snapping the optimised slot to
        mechanically-related training games.
  \item \textbf{Generalization to unseen games.}
        On a 30-game held-out set disjoint from training by name and
        token-sequence hash, we evaluate two axes of engine-level
        transfer: \emph{geometric} placement of held-out latents
        relative to training-game latents, and \emph{predictive}
        accuracy of the dynamics model on novel rule combinations.
        The results delineate where the rule-conditioned NCA
        generalizes today and where further work is needed.
\end{enumerate}
